\ifdefined\MyWebBook\else
\documentclass[../textbook.tex]{subfiles}
\begin{document}
\fi
\bookchapter{智能体运行系统}{ch:agent-runtime}

智能体运行系统负责授权、执行和记录模型提出的动作，并在中断后恢复任务。网络超时可能使执行结果暂时未知，重试可能重复产生副作用，对象状态变化后原审批也可能失效。本章以知识查询与工单处理为例，说明如何定义业务结果、持久化任务状态并处理这些故障。

\bookxref{ch:agent-architecture}讨论规划、记忆与协作；本章把这些机制落实为持久状态与执行边界。检索证据的组织见\bookxref{ch:rag-systems}，模型服务本身的连接与工作进程见\bookxref{ch:serving-architecture}。

\section{任务判定契约}

\subsection{业务执行结果分类}
设应用接收企业用户的故障描述，检索经授权的知识，读取设备状态，提出处理建议，并在允许的情况下创建或更新工单。该任务至少有三类结果：\emph{有来源的知识回答}、\emph{供人审核的操作提案}、\emph{已确认的状态变更}。三类结果应分别记录判定依据，严禁将提案完成记为业务操作完成。

任务输入应包括已认证主体、租户、用户目标、业务对象、期限与授权范围。自然语言描述可以补充目标；程序取得的身份与权限保持有效，不受文本影响。输出应携带结果类型、事实或动作引用、业务对象版本、终止原因和未完成事项。工单创建状态依据受信执行器取得的业务回执确认。本章主要符号见表\ref{tab:agent-runtime-symbols}。

\begin{table}[htbp]
\centering
\caption{智能体运行系统主要符号}
\label{tab:agent-runtime-symbols}
\begin{tabular}{ll}
\toprule
符号 & 含义\\
\midrule
$S_v,e_k$ & 版本 $v$ 的任务状态与第 $k$ 个持久事件。\\
$\delta$ & 从已记录状态与事件计算下一状态的转换函数。\\
$a,h(a),k_a$ & 动作提案、其规范化摘要与幂等键。\\
$g$ & 执行者持有的单调递增隔离代数。\\
$B,U,R$ & 任务总预算、已结算用量和未结算预留量。\\
$r_j,c_j$ & 第 $j$ 个活动的预留上界与实际成本。\\
$d,t$ & 任务截止时间和当前时间。\\
$V_o,V_p$ & 业务对象版本与权限策略版本。\\
\bottomrule
\end{tabular}
\end{table}

\subsection{运行时分层接口职责}
模型接口接收上下文并返回文本或结构化提案；工具协议描述能力发现、参数和结果；应用状态机决定是否执行、如何记录和何时结束。三层可以通过适配器连接，各自的职责在连接后仍需保持独立。

\term{模型上下文协议}{Model Context Protocol}{MCP}的2025年6月18日规范采用宿主、客户端与服务端结构，宿主管理连接及策略，协议组织上下文与能力交互\cite{mcp2025architecture}。这里引用固定版本说明分层思想，该版本只是引用基准。协议层互通不自动赋予业务数据与操作权限；新增工具服务仍需登记身份、能力、凭证范围和结果可信度。

\term{智能体间协议}{Agent2Agent Protocol}{A2A}面向彼此独立、内部实现可以不透明的智能体系统。其1.0.0规范以智能体卡片（Agent Card）声明能力与接口，以任务（Task）、消息（Message）、内容部分（Part）和制品（Artifact）表达任务状态、交互内容与产物，并支持同步、流式及长任务交互\cite{a2a2026spec100}。MCP主要连接智能体与工具、资源，A2A主要连接能够自行承担任务的智能体；二者可以组合，而应用侧的身份映射、最小授权、预算控制和业务结果核验始终由应用自己承担。\term{智能体通信协议}{Agent Communication Protocol}{ACP}曾用智能体清单（Agent Manifest）、运行（Run）和消息等对象处理相近的智能体互操作问题；该项目后来并入Linux基金会旗下A2A，原仓库已经归档\cite{acp2025archive}，属于 A2A 的演进前身，行业互操作规范统一收敛于 A2A。本节的ACP专指Agent Communication Protocol；其他资料出现相同缩写时，应核对全称与版本。

工具契约不仅包含参数类型，还要包含只读或写入性质、对象范围、前置条件、期限、重试分类、幂等支持和结果查询方式。错误也应结构化，至少区分参数错误、无权限、前置版本冲突、暂时不可用、明确未执行和执行结果未知。把所有异常折叠成“请再试一次”，最重要的执行信息就丢失了。主流智能体工具与协作协议的规范对比见表\ref{tab:tool-protocols}。

\begin{table}[htbp]
\centering
\caption{主流工具调用与协作交互协议对比}
\label{tab:tool-protocols}
\begin{tabularx}{\linewidth}{@{}p{2.2cm}p{2.8cm}p{3.2cm}X@{}}
\toprule
协议类型 & 拓扑与抽象层级 & 能力表达与交互格式 & 状态与安全边界\\
\midrule
OpenAI Tools API & 客户端--模型点对点请求响应 & JSON Schema 声明函数名、描述与参数类型；模型输出结构化字段，客户端以 \texttt{role: tool} 消息回传 & 无状态会话；依赖宿主应用注入鉴权凭证、实施白名单过滤与安全沙箱执行\\
Anthropic Tool Use & 客户端--模型流式/批量交互 & JSON Schema 声明输入模式；响应输出 \texttt{tool\_use} 块，客户端回送 \texttt{tool\_result} 块 & 支持流式交替阻断；依赖客户端侧前置格式校验与人工授权审批网关\\
MCP (Model Context Protocol) & 宿主--客户端--独立服务三层拓扑 & 基于 JSON-RPC 2.0，解耦为 Resources（只读上下文）、Prompts（工作流模版）与 Tools（动态动作） & 服务端维护生命周期与状态连接；宿主统一管理凭证下发、动态握手认证与细粒度作用域\\
A2A (Agent2Agent) & 异构独立智能体间分布式对等协作 & 以 Agent Card 声明能力清单，以 Task、Message、Part、Artifact 表达任务状态机与制品产物 & 支持同步轮询与流式事件跟踪；支持多租户身份映射、分布式预算预留与跨系统凭证隔离\\
\bottomrule
\end{tabularx}
\end{table}

\section{运行时架构}

\begin{figure}[htbp]
\centering
\begin{tikzpicture}[font=\small,node distance=7mm,
 box/.style={draw,rounded corners,align=center,text width=29mm,minimum height=12mm}]
\node[box] (api) {任务接口\\认证与限额};
\node[box,right=of api] (engine) {编排器\\状态与预算};
\node[box,right=of engine] (model) {模型适配器\\提案与解析};
\node[box,below=of engine] (db) {事务存储\\事件与发件箱};
\node[box,left=of db] (approve) {人工协作\\审批与追踪};
\node[box,right=of db] (gateway) {工具网关\\授权与隔离};
\node[box,below=of gateway] (biz) {知识与业务系统\\版本与回执};
\draw[<->] (api)--(engine);\draw[<->] (engine)--(model);
\draw[<->] (engine)--(db);\draw[<->] (approve)--(db);
\draw[->] (db)--(gateway);\draw[<->] (gateway)--(biz);
\draw[->] (gateway.north) -- (model.south);
\end{tikzpicture}
\caption[应用状态、模型建议与执行职责边界]{应用状态、模型建议与执行职责边界。应用状态负责记录业务事实，模型建议不直接产生系统副作用，外部执行负责权限核验与调用落地。}\label{fig:runtime-architecture}
\end{figure}

持久化可以从一个支持事务和唯一约束的关系数据库开始，任务事件、动作意图和发件箱在同一事务边界内提交。工作进程可水平扩展，大型证据与结果存入受控对象存储，数据库保存摘要和引用。服务拆分根据工作负载与隔离要求决定，并保留事务及状态协调机制。

\begin{table}[htbp]
\centering
\caption{智能体运行系统核心持久对象与约束}
\label{tab:agent-runtime-records}
\begin{tabularx}{\linewidth}{@{}p{.20\linewidth}p{.39\linewidth}X@{}}
\toprule
对象 & 核心字段 & 约束\\
\midrule
任务 & 租户、任务标识、状态版本、期限、预算、制品组合 & 租户与任务标识联合唯一。\\
事件 & 任务标识、序号、类型、负载摘要、发生与记录时间 & 每任务序号唯一；不以时间戳代替顺序。\\
动作 & 动作标识、工具版本、参数摘要、幂等键、状态、回执 & 同一语义动作的键唯一；尝试次数另记。\\
审批 & 动作摘要、对象版本、授权人、范围、到期时间 & 参数或对象已经变化时需重新审批。\\
发件箱 & 消息标识、动作标识、顺序、发送状态 & 与动作意图同事务写入。\\
证据引用 & 来源、版本、位置、访问范围、内容摘要 & 读取时重新满足权限与保留策略。\\
\bottomrule
\end{tabularx}
\end{table}

任务制品组合应固定模型、提示模板、工具模式、控制器和解析器版本。设备状态与权限仍随外部环境变化，执行前需要再次检查。状态快照中保存的是足够恢复的引用、版本与预算，敏感原文按引用方式处理。

\section{持久执行}

\subsection{事件重放语义}
\term{持久执行}{Durable Execution}{}使长任务在工作进程退出后能够恢复到明确状态。设初始状态为 $S_0$，则
\begin{equation}
 S_n=\delta(\delta(\cdots\delta(S_0,e_1),e_2)\cdots,e_n).
 \label{eq:runtime-fold}
\end{equation}
要使重放得到相同状态，$\delta$ 对给定输入必须确定。模型采样、外部读取和时钟是非确定输入，应先作为事件记录，再参与转换。恢复时重复调用模型会创建新的决策轨迹；跨运行时完整再现还要求固定实现、数值内核和调度条件，随机种子只是其中一项。

快照可以降低重放成本。快照记录事件序号 $k$ 与状态摘要，恢复后仅重放 $k+1$ 之后的事件。快照与事件之间必须有一致边界，否则会跳过或重复应用事件。应用升级改变转换逻辑时，要保留旧版本解释器或显式迁移事件与快照；用新逻辑无条件重放旧任务会得到与原执行不同的状态。

状态可包括已接纳、规划中、等待工具、等待审批、等待核对、已完成、已失败和已取消。后面三类为任务终态，但已取消任务仍可能保留待核对的外部动作记录：取消停止后续计划，不会自动撤回已经执行的动作。

\subsection{并发恢复版本约束}
两个工作进程可能同时读到 $S_v$。若均无条件写入下一状态，其中一个结果会覆盖另一个。\term{比较并交换}{Compare-and-Swap}{CAS}把更新限定为“当前版本仍为 $v$ 时才提交到 $v+1$”。同一旧版本最多有一个提交成功；失败者必须重新读取状态后再决定合并或重试。

租约只限制一段时间内的所有权。旧进程暂停后恢复时，租约可能早已交给新进程。\term{隔离令牌}{Fencing Token}{}用递增代数 $g$ 标记执行所有权，下游若保存已见最大代数，便可拒绝较小 $g$ 的陈旧写入。外部副作用的排他性依赖下游执行隔离令牌检查；本地租约只管理本地所有权状态。

\begin{assumption}
以下原子性论述假设事务存储对相关行提供所需隔离、唯一约束可靠，且应用确实检查提交结果。外部业务服务通过自己的幂等与结果查询接口维护执行一致性。
\end{assumption}

\section{外部动作的执行一致性}

\subsection{超时语义}
发送创建工单请求后可能发生三种情况：请求未到达、业务系统已拒绝、工单已创建但响应丢失。客户端看到的超时无法区分这些情况。若直接生成新标识再次创建，就可能得到两个工单。

\term{幂等键}{Idempotency Key}{}应绑定一个逻辑动作，在所有重试间保持稳定。服务端同时保存键、规范化参数摘要和结果；同键同参数返回已记录结果，同键异参数必须拒绝。键的作用域需要包括租户与业务操作，保留时间必须覆盖重试与核对窗口。\footnote{尝试标识用于区分网络调用，动作幂等键用于合并语义相同的调用。每次尝试生成新幂等键会失去去重作用。}

若外部服务可以原子地执行效果并登记幂等结果，则重复传输不会重复产生该效果。若其先登记“成功”再执行效果，崩溃会留下虚假完成；先执行效果再登记，崩溃则留下重复窗口。

\keyconcept{业务效果与幂等回执的原子提交，是分布式执行中避免“虚假完成”或“重复执行窗口”的核心前提。}

无法提供幂等接口的写入应优先利用业务自然键、版本条件更新或结果查询。仍无法判断时进入\term{核对}{Reconciliation}{}状态，由受控查询或人工确认结果；盲目自动重试在这个阶段只会增加重复窗口。对只读请求也要注意计费、审计和外部速率等可观察影响；这些影响计入调用成本与重试预算。

\subsection{事务发件箱}
若先提交动作意图再发送消息，提交后崩溃可能导致永远未发送；若先发送再提交，则消息可能被执行而本地不知道。\term{事务发件箱}{Transactional Outbox}{}将动作意图与待发消息放在同一数据库事务中，之后由独立分发者重试投递\cite{awsOutbox}。分发者发送成功后、更新发送状态前崩溃仍会重投，所以接收方必须去重。发件箱保障\emph{本地动作与投递意图}的一致性，接收端则负责\emph{远端业务效果}的幂等去重。

\begin{figure}[htbp]
\centering
\begin{tikzpicture}[x=2.3cm,y=.65cm,font=\small]
\foreach \x/\name in {0/编排器,1/事务库,2/分发者,3/业务服务}{
 \node at (\x,0) {\name};\draw[dashed] (\x,-.4)--(\x,-7.5);}
\draw[->] (0,-1)--node[above]{意图与发件箱}(1,-1);
\draw[->] (1,-2)--node[above]{提交确认}(0,-2);
\draw[->] (2,-3)--node[above]{读取待发}(1,-3);
\draw[->] (2,-4)--node[above]{固定幂等键}(3,-4);
\draw[->] (3,-5)--node[above]{持久回执}(2,-5);
\draw[->] (2,-6)--node[above]{记录动作结果}(1,-6);
\draw[->] (0,-7)--node[above]{按结果继续}(1,-7);
\end{tikzpicture}
\caption[幂等动作执行与业务回执机制]{幂等动作执行与业务回执机制。发送可以重复，逻辑动作身份与业务回执保持稳定。}
\end{figure}

多步业务动作部分成功时，可采用\term{补偿事务}{Compensating Transaction}{}。Saga 将长事务分成可独立提交的步骤，并为已完成步骤设计补偿\cite{garciamolina1987sagas}。补偿通过新的业务动作调整结果：已发送通知无法让接收者“未读过”，取消预订也可能产生费用。不可补偿步骤应尽量后置，并在执行前完成必要确认；补偿失败需要单独的待处理状态。

\section{异步活动控制}

\subsection{并行预算与期限传播}
若多个并行活动各自只检查当前已用成本 $U<B$，它们可能同时开始而合计超出预算。为活动 $j$ 原子预留 $r_j$，只在
\begin{equation}
 U+R+r_j\le B
 \label{eq:runtime-reserve}
\end{equation}
时接纳，其中 $R$ 为其他活动尚未结算的预留。完成后执行 $U\leftarrow U+c_j$、$R\leftarrow R-r_j$。若 $0\le c_j\le r_j$ 且提交原子，则 $U+R\le B$ 在归纳意义下保持成立。若计费方允许实际成本超过预留，上述保证不再成立，必须限制生成与工具上限或设置额外风险准备。

\begin{example}[并行调用的预算预留]
预算100单位、已用40、已预留20，新活动预留30可接纳，总占用90；另一活动再申请20必须等待或拒绝。前一活动实际消耗24后，状态变为 $U=64,R=20$，只剩16单位可继续预留。
\end{example}

期限同样应从父任务向子活动传播。子活动可用时间不超过 $d-t$，并预留持久化与清理时间。重试次数、模型词元、工具次数、费用与墙钟期限是不同约束，任何一个触达上限都应产生明确终止或等待决策。

\subsection{取消与完成的竞争}
取消请求首先形成持久事件，然后停止新动作分派，并向可取消活动发送信号。若动作已经提交，取消只能中止等待，效果留在外部系统中。完成回执与取消事件竞争时，状态机应保留两项事实：用户已要求停止，以及某外部动作最终是否完成。晚到回执仍应写入动作记录，以保存实际业务效果。

重试应限于确定可重试的错误，并受剩余期限和幂等能力约束。指数退避可以减少持续冲击，随机扰动降低同时恢复的同步峰值；参数错误、拒绝授权与已确认的业务冲突都要在重试之外处理。异步回调需要认证、动作关联和去重；外部传来的任务标识只是待核对线索，据此直接修改状态会绕过身份校验。

\section{工具执行授权}

数据库工具优先暴露受限业务查询，任意高权限查询接口只保留给运维通道；代码执行工具需要进程或容器隔离、资源限制、文件与网络范围，凭证由网关注入，不进入模型上下文；业务写入工具应支持对象版本和结果查询。检索权限则在取回内容前或受信过滤层实施，生成后删除敏感词补偿不了缺失的授权。

\subsection{智能体的权限模型}
传统应用与智能体应用都要核对主体、对象和操作，也都可以使用角色、属性、委派及逐次鉴权。区别在于智能体的下一步动作由模型根据运行中的观测动态提出，提案参数还可能受外部网页或工具结果影响；用户授权“整理工单”并未预先批准模型随后提出的任意查询、文件读取或状态变更。Harness 因而要把模型输出视为待审核的动作候选，将已认证用户、任务授权范围、执行者身份、目标对象、具体操作及当前状态一并送交受信策略检查。工具出现在模型可选列表里，只表明模型能够提出调用，调用能否落地仍由执行层决定。

权限可按三个相接的范围理解：用户和业务系统允许访问哪些对象；当前任务获准使用其中哪些工具与操作；执行环境实际能读写哪些文件、访问哪些网络目标及取得哪些凭据。每次动作须同时满足这三类限制，审批还须绑定具体动作与对象版本，而非笼统绑定一次会话。例如用户允许智能体读取工单123并起草回复，读取其他租户的工单、发送邮件或关闭工单仍需分别取得相应授权。模型建议的工具参数即使符合模式，也必须核对真实对象的归属和当前业务状态；已经取回的资料不能因模型再次引用就扩大披露范围。

委派子智能体时，应从父任务的剩余权限与预算中划定更小的范围，并保留父子任务身份供审计；子智能体的角色描述和项目规则不授予额外权限。工具网关负责检查业务对象和动作，沙箱约束命令与代码在运行环境中的可达资源，目标业务系统继续执行自己的鉴权。三处检查分别面对模型提案、进程行为和最终业务写入，不能仅依赖其中一处。

\subsection{Harness 的执行沙箱}
智能体在处理任务时可能提出运行脚本、调用命令行工具或执行外部代码。模型能检查提案文本，却无法据此穷尽程序运行后的文件访问、网络请求和资源消耗；直接继承宿主进程权限，还可能让一次错误命令触及其他任务的数据。\term{执行沙箱}{Execution Sandbox}{}为这类动作提供受约束的运行环境，把可访问的文件、网络、计算资源和执行时间限制在任务允许的范围内。沙箱约束的是程序实际运行时的能力，不能靠提示词中的“不要访问”来实现。

一种运行架构将管理任务的控制层与实际执行命令的环境分开：Harness 保存任务身份、授权和预算，为任务分配隔离工作区及执行期限；工具网关先核对动作与审批，再把允许的命令交给沙箱内的执行服务。运行环境通过文件挂载、资源配额和网络策略限制进程行为，执行服务仅通过受控接口返回输出、退出状态和产物。于是，网关负责判断“这次动作能否发起”，沙箱负责约束“动作发起后能做什么”。沙箱可以用容器或虚拟机等方式承载，选用何种隔离方式取决于工作负载与威胁边界。业务 API 的写入则还要接受目标系统的鉴权与回执核验；并非每个工具都需要作为沙箱内的程序运行。

以代码修复任务为例，模型提出运行测试，Harness 将指定项目文件提供给隔离工作区，只开放测试所需的临时写入位置、时间和网络范围。测试脚本即使试图读取另一项目的文件，也应被运行环境拒绝；测试失败不会赋予模型更大的权限。测试完成后，Harness 取回差异与日志，记录执行结果并清理环境；任务中断后恢复时，需重新确认沙箱状态和授权。需要访问外部服务的动作由网关按范围提供短期凭据，不把长期密钥放进模型上下文或可读取的沙箱文件。子智能体也只能在父任务获准范围内获得工作区与工具能力。

验证沙箱要检查实际可达路径，而不只检查配置名称。容器若挂载敏感目录或高权限套接字，进程仍可使用这些入口；只限制域名解析也不能约束直接连接已知 IP 的请求。验收时应同时检查越界读写、网络访问、资源超限与任务间残留是否受阻，以及正常测试能否完成。拒绝、超时和环境失败返回不同状态供编排器处理；命令成功退出也只确认该次执行结果，外部业务状态仍以目标系统回执为准。

\term{检查与使用时差}{Time-of-Check to Time-of-Use}{TOCTOU}发生在检查之后、执行之前对象或权限发生变化。审批应绑定 $h(a)$、$V_o$、授权范围和到期时间。执行前再次读取对象版本与当前权限；不一致时重新规划或重新审批。批准“把工单123分派给甲组”只覆盖该动作摘要，后来被模型改成的“关闭工单123”需要另行审批。

\begin{warning}[审批与执行时差（TOCTOU）失效]
人工审批必须强绑定动作摘要与当前业务对象版本；若在审批等待期内目标对象状态或策略版本已发生变更，原审批决议自动失效，执行器严禁依据过期版本直接提交操作。
\end{warning}

人工协作是一种明确的等待状态，保存展示给人的提案、证据和选择。审批通过后恢复的是同一个动作身份；拒绝应被记录为约束或终止，模型换一种措辞重发相同动作时按同一被拒记录处理。对于本来已获授权且参数未变的低风险重复执行，重复审批会增加时延而不增加安全性。

\section{智能体运行平台}
\optional
\label{sec:harness-platforms}
本节将具体项目作为架构案例。以下产品事实依据2026年9月8日查阅的一手资料；功能会随版本变化，选型时还需锁定所采用的发布版或提交。表\ref{tab:agent-runtime-harness-platforms}概括各平台的组织方式，平台之间可通过协议组合。

\begin{table}[htbp]
\centering
\caption[典型智能体运行时架构关注点对比]{典型智能体运行时架构关注点对比。展示四种具体实现的架构侧重；表格不表示性能排名或生产验收结果。}\label{tab:agent-runtime-harness-platforms}
\begin{tabularx}{\linewidth}{@{}p{2.5cm}XX@{}}
\toprule 项目 & 可观察的组织方式 & 应审查的工程问题\\\midrule
Codex & App Server 将交互客户端接入统一的智能体执行能力 & 线程与轮次身份、事件流、审批与断线恢复\\
DeepSeek Harness & 以 Cordis 插件组织模型、工具、存储与运行循环 & 插件依赖、替换边界、版本兼容与日志语义\\
OpenClaw & 常驻 Gateway 连接消息渠道、控制客户端与设备节点 & 消息路由、设备身份、连接恢复和授权主体\\
Hermes Agent & 跨会话记忆、技能维护、消息网关与调度 & 经验写入质量、技能版本、长期任务及用户隔离\\
\bottomrule
\end{tabularx}
\end{table}

\subsection{Codex 与 DeepSeek Harness}
Codex 的 App Server 将智能体执行能力提供给客户端，交互协议包含线程、轮次、进度事件和审批请求。OpenAI 的架构文章描述了基于 JSON-RPC 的集成方式，并说明界面与执行状态分离的意义：客户端断开后，服务端的任务事实继续保存。这种设计使多个交互入口共享执行核心，实际工具与权限由部署配置确定。\cite{openai2026codexharness}

DeepSeek Harness 的官方说明以 Cordis 插件系统为核心，模型、工具、技能、会话、沙箱、存储、循环、调度与界面均可通过插件组合；会话记录采用追加式日志。资料查阅时它仍处于开发者预览阶段。\cite{deepseek2026harness}插件化使运行循环也可替换，替换时需维持组件之间的语义兼容。例如，替换存储插件时必须确认事务与恢复语义；替换工具插件时必须保持动作身份及结果契约。采用该项目时应锁定版本，并检查接口稳定性与部署要求。

\subsection{OpenClaw 和 Hermes Agent}
OpenClaw 的官方架构以常驻 Gateway 统一连接消息渠道，控制客户端和设备节点通过 WebSocket 接入；设备能力与身份是显式协议字段。文档还指出事件出现缺口时客户端需要刷新状态。\cite{openclaw2026architecture}因此网关既是渠道适配点，也是会话路由与设备信任的汇合点。设备操作权限根据消息主体与设备授权关系校验。

Hermes Agent 是 Nous Research 的智能体项目，区别于同名模型系列。其官方仓库描述跨会话记忆、会话检索、可维护技能、消息网关和调度能力\cite{nous2026hermesagent}。该系统的“从经验学习”指外部结构化记忆与可复用技能文件的动态维护，并不改变底层模型参数权重。任务成功后保存的技能，在后续调用时仍需校验适用前提、来源版本与失败约束；环境变化后须重新核验有效性，并支持失效经验的显式撤回。

这类个人智能体还面临渠道身份与执行身份不一致的问题。例如，群聊发言者、消息渠道账号、操作系统账号和业务系统账号是不同主体。可靠实现应明确从哪一个主体派生权限，在哪一层检查对象范围，以及一次会话能否访问另一会话的记忆。各渠道按相同的主体映射与对象授权规则执行。

\subsection{运行平台的选用}
比较产品时，应先给出任务闭环，再逐项映射组件。代码修改需要工作区隔离、差异审查与测试证据；持续个人助理需要身份路由、可取消调度和记忆治理。若个人助理委派编码 Harness，应记录父子任务身份、预算、工作区、允许工具和最终产物，下游自然语言“完成”只是待验证主张。集成应通过明确的工具或客户端协议进行；拼接两个聊天记录来推断执行顺序会丢失身份与顺序信息。

\section{项目上下文管理}

智能体平台常用 Markdown 或相近的文本文件保存长期有效的项目上下文。文件名本身不决定运行语义：Harness 必须先发现文件，再按作用域、优先级和上下文预算决定注入、检索或忽略。表\ref{tab:agent-project-context-files}按职责归纳常见形式；同一行中的名称是不同平台解决相近问题的实现，不要求一个项目同时创建全部文件。

\begin{table}[htbp]
\centering
\caption{智能体项目常见上下文与知识文件}
\label{tab:agent-project-context-files}
\begin{tabularx}{\linewidth}{@{}p{2.1cm}p{4.1cm}XX@{}}
\toprule
职责 & 常见文件或目录 & 适合保存的内容 & 加载与治理要点\\
\midrule
项目指令 & \texttt{AGENTS.md}、\texttt{CLAUDE.md}、\texttt{GEMINI.md}、\path{.github/copilot-instructions.md} & 构建命令、目录边界、代码规范、验收方式和安全约束 & 由兼容的 Harness 自动发现；根目录与子目录可以形成作用域链，应保持简短并显式处理冲突\\
身份与用户模型 & \texttt{SOUL.md}、\texttt{IDENTITY.md}、\texttt{USER.md} & 角色定位、表达风格、交互边界及稳定的用户偏好 & 属于部分个人智能体平台的启动上下文，不是通用编码代理标准；私人信息应限制访问和传播范围\\
持久记忆 & 平台管理的本地记忆目录、\texttt{MEMORY.md}、\path{memory/YYYY-MM-DD.md} & 已确认的长期事实、决策摘要、任务经历与后续线索 & 文件名和加载方式取决于平台；写入时保留来源、时间、对象范围和失效条件\\
设计事实源 & \texttt{DESIGN.md}、\texttt{ARCHITECTURE.md}、\path{docs/design-docs/*.md}、ADR & 架构边界、接口契约、设计理由、替代方案与已接受决策 & 通常是供人和智能体共同读取的工程文档；由项目指令提供显式索引规则，平台默认不主动扫描未引用文档\\
技能与专用代理 & \texttt{SKILL.md}、\path{.agents/skills/*}、\path{.github/agents/*.agent.md} & 可复用流程、适用条件、工具说明、模板和专门角色 & 一般按任务相关性或显式请求载入；技能说明不改变工具权限，目录位置和元数据格式取决于平台\\
\bottomrule
\end{tabularx}
\end{table}

\subsection{项目规则的作用域与信任边界}
\term{项目规则}{Project Rules}{}为编码智能体提供持续适用的工作约定，例如目录职责、构建与测试命令、代码风格、数据处理限制和完成标准。规则可以位于项目根目录、子目录或按路径匹配的文件中；Harness 负责发现、装配及冲突处理。以已认证任务和平台策略为准确定规则的适用范围，再加载任务所涉及目录的规则；子智能体接收自身工作范围内适用的规则，并继承父任务的权限与预算上限。具体文件格式、加载时机和同级冲突次序依所选平台的版本而定。

规则文件宜只保存可重复执行的约束与入口，把大型设计说明、单次任务计划和可变事实放在各自的文档或状态存储中。若根目录要求运行测试，而局部目录说明必须先生成代码，执行时需同时满足两项约束；若两条适用规则矛盾，应按平台支持的优先级和明确的项目决策解决，无法消解时暂停受影响动作并请求澄清。外部网页、仓库依赖和工具输出中的命令式文字属于待处理数据，不因出现在上下文中就成为项目规则。

规则决定智能体应如何提出和检查动作，真正的读写权限、审批与沙箱范围由受信执行层强制实施。规则声称“允许访问”也不能扩大已经配置的文件、网络或业务对象权限；子智能体和技能继承的仍是同一执行边界。常见项目指令文件及其平台加载差异见本节后文。

\subsection{规格驱动开发}
\term{规格驱动开发}{Spec-Driven Development}{SDD}把可审查的规格作为编码智能体实施功能的起点。规格描述用户场景、功能需求、边界条件和验收标准；计划把规格映射为架构、接口与技术约束；任务清单形成可执行且可追踪的工作单元；实现完成后，再以测试、静态检查、构建结果和运行观测核对规格。GitHub Spec Kit 将这一闭环概括为 Specify、Plan、Tasks、Implement 和 Converge，并在需求含糊、制品不一致等位置加入澄清与分析门禁\cite{githubWhatSpecdrivenDevelopment2026}。

SDD 中的制品承担不同职责：项目指令规定长期工程规则，功能规格定义期望行为与验收条件，设计文档和 ADR 记录实现边界及决策理由，代码与验证证据说明当前实现状态。需求变化时，应更新受影响的规格、计划、任务和验收证据，并检查它们与代码的一致性。这个过程为智能体提供结构化上下文和完成判据；规格的质量、实现正确性与生产可用性仍由评审和验证活动共同保证。

相关工具覆盖的工程层次各有侧重。OpenSpec 围绕提案、规格、设计、任务、实施和归档组织可版本化的变更生命周期；GSD 将元提示、上下文工程与 SDD 结合，以分阶段规划和执行控制长任务的上下文；Superpowers 把需求澄清、设计审查、实现计划、测试驱动开发、子智能体执行和代码评审封装为可组合技能\cite{openspeccontributorsOpenSpec2026,tachesGSDDocumentation2026,vincentSuperpowers2026}。选型时可分别评估规格制品、上下文装配、执行编排和质量门禁；组合使用时，应指定规格、计划和任务状态各自唯一的事实源与门禁次序，减少重复生成和状态漂移。

以2026年9月22日可查阅的官方文档为例，Codex 从全局目录、项目根目录到当前工作目录构造 \texttt{AGENTS.md} 指令链，较近目录的规则覆盖较远规则；Claude Code 与 Gemini CLI 分别使用分层的 \texttt{CLAUDE.md} 和 \texttt{GEMINI.md}；GitHub Copilot CLI 还支持仓库级及按路径匹配的指令文件\cite{openaiCustomInstructionsAGENTSmd,anthropicManageClaudeMemory,googleProvideContextGEMINImd,githubAddingCustomInstructions}。OpenClaw 则把操作规则、人格、用户模型与长期记忆拆分到不同文件，并为不同文件设置不同的加载范围\cite{openclawAgentWorkspace}。这些名称说明产品约定，迁移平台时仍需核对实际发现路径、覆盖次序、导入语法、大小限制和子智能体继承规则。

\subsection{Codex 的本地记忆流水线}

Codex 的本地客户端使用独立于 ChatGPT Web 记忆的本地存储。启用记忆后，系统从符合条件的历史任务中抽取可复用信息，在后台生成摘要、持久条目、近期输入和支撑证据；正在进行或持续时间很短的任务可以不进入抽取，后台处理也不要求在任务结束时立即完成。默认存储位于 \path{~/.codex/memories/}，其中内容属于可检查的生成状态，不应作为人工维护的唯一事实源。会话级控制分别决定当前任务能否读取既有记忆，以及当前任务能否贡献给后续记忆；全局设置控制功能启用和抽取、归并行为\cite{openaiCodexMemories}。

图\ref{fig:codex-local-memory-pipeline}把稳定指令路径与可召回记忆路径分开。\texttt{AGENTS.md} 按目录作用域形成基础指令 $I_t$；历史任务先经过抽取和归并，任务发生时再按相关性形成 $M_t$。记忆内容按当前任务的相关性召回；持续适用的构建、权限和验收规则则写入相应作用域的指令文件\cite{openaiCustomInstructionsAGENTSmd,openaiCodexMemories}。

\begin{figure}[htbp]
\centering
\includegraphics[width=0.98\linewidth]{\subfix{../../figures/theory/codex-local-memory-pipeline.pdf}}
\caption[记忆与指令进入上下文的双重路径]{记忆与指令进入上下文的双重路径。记忆先形成可检索状态，再按当前任务召回；会话和全局控制分别约束读取与贡献。}\label{fig:codex-local-memory-pipeline}
\end{figure}

人工纠正适合建模为带来源、时间和作用域的新观测，再由归并过程建立新增、替代或失效关系。追加式更正能够保留修订链，生成索引和摘要则可以重新构建；直接覆盖派生摘要会丢失原始依据。具体目录和文件名属于实现约定，系统设计依赖的是可追溯的更正记录、确定的归并语义和可撤销的发布过程。

本地保存也需要数据治理。自动脱敏降低生成字段包含秘密的概率，仍应在共享 Codex 主目录或记忆制品前检查内容；密码、令牌和私人原始记录不进入记忆。若记忆来自网页、消息或计算机操作历史，还应把来源内容视为不受信输入，在召回后继续执行权限过滤和提示注入防护。

\subsection{注入、检索与事实源}
设自动注入的适用指令为 $I_t$，从设计文档和技能中按需取得的知识为 $K_t$，从记忆系统召回的记录为 $M_t$，当前工具观测为 $O_t$，则送入模型的任务上下文可以写为
\begin{equation}
 c_t=\operatorname{Assemble}(q_t,I_t,K_t,M_t,O_t;B_{\mathrm{ctx}}),
\end{equation}
其中 $q_t$ 是当前请求，$B_{\mathrm{ctx}}$ 是上下文预算。项目指令通常进入每次适用任务的基础上下文；大型设计文档、历史记录和技能资源更适合按需检索。把全部资料复制进一个规则文件会挤占任务、代码和证据的空间，也使局部规则难以独立演进。较稳健的组织方式是让项目指令充当地图，把架构与产品事实保存在结构化文档中，并用明确链接指向它们\cite{lopopolo2026harnessengineering}。

四类内容还应采用不同的更新条件。项目指令随工程规范变更而评审；设计文档随架构决策和接口版本更新；记忆条目随新证据进行替代、降权或失效；技能按输入输出契约和工具版本发布。\texttt{.agents/}、\texttt{.claude/}、\texttt{.cursor/}等隐藏目录只是平台命名空间，目录中的规则、代理、技能或设置仍应按各自语义治理。密钥、令牌和私人原始记录不进入版本化指令或设计文档；Harness 通过凭证存储、运行时注入和访问控制向工具提供必要权限。

\section{计算机交互控制}

\label{sec:computer-use}
\term{计算机操作}{Computer Use}{}指智能体通过屏幕、界面结构及输入动作操作软件的能力。它可覆盖浏览器与原生应用，其机制是观察界面、提出动作、由执行器实施，再读取新状态。模型输出的一次点击本身不驱动操作系统；执行环境必须实现相应动作并回传观测。OpenAI 的计算机操作接口文档也采用这种模型与执行环境分工。\cite{openai2026computeruse}

\subsection{界面坐标映射}
记界面观测为 $o_t=(I_t,D_t,A_t,v_t)$：$I_t$ 是截图，$D_t$ 是可获得的文档对象模型，$A_t$ 是可获得的无障碍树，$v_t$ 是页面或窗口状态标识。各环境提供哪些分量取决于具体实现。截图适用于视觉布局与画布，但坐标依赖窗口、缩放和滚动；文档结构适合语义定位，像素内容不在其中；无障碍树暴露角色、名称和状态，其完整性取决于应用实现。

\term{动作定位}{Grounding}{}把“打开保存菜单”这样的意图映射到当前界面的具体目标。动作应绑定窗口或标签页、目标描述及观测版本；只保存一个坐标容易在页面变化后点击错误对象。若在原始截图 $(W_s,H_s)$ 上定位到 $(x_s,y_s)$，截图对应视口区域的原点为 $(x_0,y_0)$、大小为 $(W_v,H_v)$，且只有轴对齐缩放，则
\begin{equation}
 \begin{aligned}
  x_v &= x_0+x_s\frac{W_v}{W_s}, \\
  y_v &= y_0+y_s\frac{H_v}{H_s}.
 \end{aligned}
\end{equation}
各量必须使用约定的像素或逻辑坐标单位。只有知道截图尺寸与视口映射后才能换算；缺乏坐标约定时，设备像素比再次乘入会得到错误的动作坐标。裁剪、浏览器工具栏偏移、iframe 坐标及页面滚动会改变映射，应重新获取对应几何关系。

\subsection{动作结果验收}
执行链必须区分三件事：\emph{输入动作已派发}、\emph{界面出现预期变化}、\emph{业务后置条件成立}。点击提交后按钮变灰只支持第一或第二项；确认对象已经创建，通常需要成功回执、对象标识或结果查询。对于回执未知的提交，应先查询结果，再决定是否重试。没有幂等语义的界面重放整段动作会产生重复业务效果。

设任务有 $n$ 步，第 $i$ 步在此前步骤均正确的条件下成功的概率为 $p_i$，则由条件概率链式法则
\begin{equation}
 P(\text{全程成功})=\prod_{i=1}^{n}p_i.
\end{equation}
此式不要求步骤独立，但 $p_i$ 必须是上述条件概率；不同测试的边际点击准确率直接相乘会得到偏乐观的估计。长流程会累积失误，因此应在有业务意义的阶段读取后置条件，及时停止错误传播。阶段检查点保存进度，已提交的外部动作通过幂等查询与补偿处理。

\subsection{界面操作权限}
页面中的文字是待处理数据，出现在屏幕上不改变它的身份。截图、DOM 和无障碍树都可能承载提示注入；结构化读取只能改善定位，恶意内容仍需另外隔离。执行器应绑定用户已授权的目标、对象和动作范围，隔离会话数据，并在界面变化时重新验证目标。截图与轨迹也可能包含隐私数据，需要控制保存范围、访问权限和保留时间。

\section{浏览器控制协议}
\optional
\label{sec:browser-control}
\term{Chrome 开发者工具协议}{Chrome DevTools Protocol}{CDP}用于检测、调试和控制 Chromium 系浏览器。它按 DOM、Network、Runtime、Page 等域组织命令与事件；协议消息使用结构化 JSON\cite{chromium2026cdp}。CDP 提供浏览器内核调试与自动化控制原语，而计算机操作面向端到端的用户任务动作抽象。智能体既可以通过 CDP 读取 DOM、网络与运行时状态，也可以基于屏幕截图与模拟输入事件驱动页面交互。

\subsection{浏览器会话管理}
浏览器包含多个 target，例如页面与工作线程；控制器需要先选定目标，再将命令与相应会话关联。请求编号用于匹配响应，事件用于接收异步变化。一次导航可能使先前保存的节点句柄或执行上下文失效，弹窗可能创建新的目标，因此“当前活动页”只是临时指向，不足以充当稳定的业务对象身份。控制器应保存目标身份并在导航、断线或状态缺口后重新读取。

在工具层可以定义“定位当前页面中名称为保存的按钮并确认可操作”，由适配器负责浏览器细节。Playwright 提供这类更高层的浏览器自动化接口，也支持通过 CDP 连接 Chromium；其文档明确提示该连接方式相较原生 Playwright 协议连接具有较低保真度。\cite{playwright2026browsertype}连接路径应按任务所需能力选择。\term{WebDriver 双向协议}{WebDriver BiDi}{}是另一条标准化浏览器自动化路径，提供双向通信和事件机制；具体浏览器的实现覆盖仍需核对。\cite{w3c2026webdriverbidi}

\begin{table}[htbp]
\centering
\caption[操作接口选型依据与评估维度]{操作接口选型依据与评估维度。选择依据对象语义、可验证性和权限；能否点击成功只是其中一项。}\label{tab:agent-runtime-browser-interfaces}
\begin{tabularx}{\linewidth}{@{}p{2.3cm}XX@{}}
\toprule 接入方式 & 适合表达的对象 & 主要限制\\\midrule
业务 API & 业务对象、事务、明确回执 & 需要存在可用且获授权的接口\\
浏览器自动化库 & 定位器、页面、等待条件 & 依赖应用结构和驱动兼容性\\
CDP & 浏览器目标、文档、网络与运行时 & 主要面向 Chromium，需处理版本和会话\\
无障碍接口 & 控件角色、名称与状态 & 应用暴露的语义可能不完整\\
截图与输入 & 可见像素与坐标动作 & 定位误差、几何变化与回执歧义\\
\bottomrule
\end{tabularx}
\end{table}

\subsection{兼容性和执行边界}
CDP 的 tip-of-tree 定义会变化，官方不承诺其向后兼容；应锁定浏览器与协议适配版本，检查所需命令是否受支持。\footnote{CDP 文档中的 stable 1.3 是历史冻结子集，覆盖范围小于现代浏览器的全部能力；tip-of-tree 跟随开发版本变化，使用时需固定具体协议版本并核查兼容性。}远程调试连接具有广泛的会话访问能力，应限制其可达范围，并与用户日常浏览配置隔离。使用浏览器已有登录态必须明确授权范围；获得调试连接端口并不等同于获得全部会话标签页与个人数据的访问授权。

MCP 等工具接入协议可以把浏览器能力暴露给模型，实际适配器再调用 CDP 或其他驱动。运行系统负责端到端幂等、认证与业务成功判定，工具协议负责能力接入。测试除正常操作外还应覆盖导航中断、目标关闭、登录过期、遮挡、重复按钮、回执丢失和恶意页面指令，并分别记录定位、执行与业务验证失败。

\section{执行恢复}

业务操作的完成状态由目标系统的回执确认。运行系统应将回执与动作身份、对象版本绑定并持久化，再据此生成用户答复。发生中断时，恢复过程先核对已提交的动作和回执，随后决定继续投递、查询结果或重建答复。

\begin{example}[从执行结果生成业务回执]


考虑用户要求“根据设备A的告警及维护规定建立待处理工单”。执行链路如下。
\begin{enumerate}
\item 接口建立任务，固定租户、主体、设备标识、期限和应用制品版本；以请求幂等键处理用户重复提交。
\item 检索器按主体权限取得维护文档版本与位置；设备查询工具返回状态及对象版本17。二者均保存来源与时间，文档规范与实时遥测数据分别建立独立版本与时间戳记录。
\item 模型产生工单草案，包含设备、告警摘要、支持证据和建议优先级。解析器检查结构，业务规则检查必填信息与取值范围；缺失信息转入等待用户补充。
\item 若操作在授权范围内，执行器绑定参数摘要与设备版本17；若需审批，持久化该提案并等待。期间对象变化则重新评估。
\item 在事务中保存动作意图、预算预留及发件箱。工具网关再次授权，以稳定动作键调用工单系统。
\item 工单系统返回工单标识与版本。运行系统记录回执后才生成“已创建”的答复；若超时，则按键查询或进入待核对，答复必须说明状态尚未确认。
\item 后续更新工单采用其版本条件。任务结束释放未用预留，保留可追溯事件与按保留策略管理的证据引用。
\end{enumerate}

\begin{figure}[htbp]
\centering
\begin{tikzpicture}[x=1cm,y=1cm,>=stealth,font=\small]
\foreach \name/\xx/\title in {o/0/编排器,d/2.5/事务存储,g/5/工具网关,b/7.5/业务系统,u/10/用户}{
 \node[draw,align=center,text width=17mm,minimum height=8mm] (\name) at (\xx,0){\title};
 \draw[dashed] (\xx,-.45)--(\xx,-7.7);
}
\node[align=left,font=\scriptsize,anchor=west] at (-.4,-.8){前提：结构、权限、前置条件与所需批准已满足；否则不进入提交。};
\draw[->] (0,-1.4)--node[above,font=\scriptsize]{条件事务提交}(2.5,-1.4);
\node[align=center,font=\scriptsize] at (2.5,-2){动作身份、预算预留\\意图事件与发件箱};
\draw[->] (2.5,-2.7)--node[above,font=\scriptsize]{固定键与参数摘要}(5,-2.7);
\node[font=\scriptsize] at (5,-3.2){当前授权重查};
\draw[->] (5,-3.7)--node[above,font=\scriptsize]{同键提交}(7.5,-3.7);
\draw[->] (7.5,-4.3)--node[above,font=\scriptsize]{业务回执}(5,-4.3);
\draw[->] (5,-4.9)--node[above,font=\scriptsize]{持久化结果并结算}(2.5,-4.9);
\draw[->] (2.5,-5.5)--node[above,font=\scriptsize]{确认后的动作状态}(0,-5.5);
\draw[->] (0,-6.1)--node[above,font=\scriptsize]{已创建＋工单标识；由回执生成答复}(10,-6.1);
\node[align=center,font=\scriptsize] at (5,-7){若业务回执丢失：先登记结果未知，再按原键查询业务系统；\\确认结果后复用同一结果持久化路径，不重新创建逻辑动作。};
\end{tikzpicture}
\caption[持久化工单创建执行时序]{持久化工单创建执行时序。当前授权重查失败时终止发送并记录拒绝；业务回执缺失时只能给出“状态尚未确认”的答复。}\label{fig:runtime-tool-sequence}
\end{figure}

图\ref{fig:runtime-tool-sequence}区分本地事务提交、远端业务提交和最终回答三个时刻。发件箱只保证发送意图不会被本地故障静默丢弃；远端幂等、结果查询与重复回执处理共同约束恢复行为。
\end{example}

\subsection{故障恢复路径}
\begin{table}[htbp]
\centering
\caption{不同故障位置事实与恢复动作}
\label{tab:agent-runtime-recovery}
\begin{tabularx}{\linewidth}{@{}p{.28\linewidth}p{.32\linewidth}X@{}}
\toprule
故障位置 & 可确认事实 & 恢复动作\\
\midrule
意图事务提交前 & 无持久动作承诺 & 从最近事件恢复，不假定已经发送。\\
提交后、发送前 & 发件箱存在 & 分发者继续投递。\\
远端完成、响应丢失 & 本地结果未知 & 以相同键查询或重试，禁止新建逻辑动作。\\
结果记录后、答复前 & 业务回执已持久化 & 由回执重建答复，不再次执行。\\
审批后对象变更 & 原审批基于旧版本 & 再检查、重规划，必要时重新审批。\\
取消后收到回执 & 取消与外部完成均发生 & 记录真实效果，按授权决定是否补偿。\\
\bottomrule
\end{tabularx}
\end{table}

\begin{algorithm}[受控动作提交与恢复]

\begin{enumerate}
\item 从任务版本读取允许工具、剩余预算、期限和终止状态，解析提案并生成规范化参数摘要。
\item 检查当前权限、业务前置条件与审批绑定；不满足则记录明确状态，不调用外部写入。
\item 以任务版本条件，在同一事务中预留预算、建立动作身份、追加意图事件和发件箱记录。
\item 分发者使用固定幂等键执行；把明确失败、成功回执和结果未知分别记录。并发记录通过动作版本与唯一约束去重。
\item 对未知结果使用受信查询核对；只有满足幂等契约才自动重试写入。所有尝试共享原动作身份。
\item 结算实际成本并追加结果事件，驱动状态转换；恢复时读取既有结果，不重复已完成副作用。
\end{enumerate}
\end{algorithm}

\subsection{部署拓扑和故障隔离}
任务接口可以无状态复制，工作进程按模型活动与工具活动分池，以避免慢工具占满模型执行槽。各副本通过事务存储竞争任务版本；进程内存中的任务状态在副本重启后即失效。知识索引按权限域访问，业务系统仅对工具网关开放凭证范围，模型服务没有直连业务写入的权限。

\begin{figure}[htbp]
\centering
\begin{tikzpicture}[font=\small,node distance=7mm,
 box/.style={draw,align=center,text width=30mm,minimum height=12mm}]
\node[box] (edge) {接入区\\接口副本与认证};
\node[box,right=of edge] (work) {执行区\\编排与活动副本};
\node[box,right=of work] (model) {模型区\\推理服务池};
\node[box,below=of edge] (store) {数据区\\事务库与对象存储};
\node[box,below=of work] (gate) {隔离区\\工具网关与沙箱};
\node[box,below=of model] (business) {业务区\\知识库与工单系统};
\draw[->] (edge)--(work);\draw[<->] (work)--(model);
\draw[<->] (work)--(store);\draw[<->] (work)--(gate);
\draw[<->] (gate)--(business);
\end{tikzpicture}
\caption[推理服务与业务写入的隔离拓扑]{推理服务与业务写入的隔离拓扑。以网络与凭证范围隔离模型推理和业务写入。部署区表示职责边界，不规定必须使用六套集群。}\label{fig:runtime-deployment}
\end{figure}

图\ref{fig:runtime-deployment}中的边界由网络与凭证共同落实，服务命名本身不提供隔离。事务存储需要满足明确的恢复点目标与恢复时间目标；备份恢复还须与外部已完成动作核对。若数据库恢复到了动作执行前的旧时间点，简单重放发件箱可能再次发送已执行动作，因此幂等记录的保留窗口应覆盖灾难恢复范围。跨区恢复除了验证服务启动，还要验证任务身份、动作键和业务回执是否继续一致。

\section{运行系统的质量保障}

\subsection{端到端执行可观测性}
一次任务可跨越检索、模型调用、人工审批与多次工具尝试。用任务标识关联这些步骤，以动作标识聚合针对同一业务动作的尝试；每次调用记录开始、结束、结果类别、耗时与成本，并关联事件序号、模型修订及工具版本；相应步骤再记录证据引用、终止原因或业务回执。跨组件追踪展示步骤的时间关系，持久事件与业务回执提供执行事实。异步任务恢复后沿用任务标识，重新调用工具时保留原动作标识并生成新的尝试标识，才能区分续跑和重复执行。指标、日志与追踪的一般分工见\bookxref{ch:serving-architecture}[中的“服务遥测”一节]。

例如创建工单时工具请求超时，追踪指出耗时发生在工具网关还是业务系统，却无法单独判定工单是否已创建。运行时应按动作标识查询业务回执或对象状态，将结果归为已确认完成、明确未执行或待核对；待核对仍占用未完成任务的统计分母。排查时按任务标识串起上下文版本、动作提案、准入决定、尝试结果和最终业务状态，便可区分模型选错工具、授权拒绝、传输超时和业务写入冲突。

监测分别统计任务完成率、工具实际完成率、拒绝与待核对比例、人工介入率，以及任务级端到端时延和成本；按任务类型、制品修订与结果类别分组，报告观察窗口和样本量。重试成本与等待审批的时间计入任务总量，失败和取消任务保留在分母中。指标标签采用有限维度，任务和尝试的唯一标识放入受控追踪或日志。记录默认保存必要标识、版本、决策类别和受控负载引用；提示原文、检索材料、工具响应及截图按敏感数据管理，不无差别采集模型内部思维过程。

\subsection{执行轨迹的判定边界}
\term{轨迹回放}{Trajectory Replay}{}使用已记录模型结果与工具结果验证控制器状态转换；回放用于检查分支、预算与恢复逻辑。真实集成验证需要在受控业务环境中核对接口语义和最终状态。故障注入应覆盖发送前后崩溃、重复回调、过期租约、参数冲突、审批撤销、取消竞争与预算并发。

评估应分别测量参数有效、授权判定、工具实际完成和最终任务完成。对串行必要环节 $E_1,\ldots,E_m$，准确关系是
\begin{equation}
 \Pr\left(\bigcap_{j=1}^mE_j\right)
 =\Pr(E_1)\prod_{j=2}^m\Pr(E_j\mid E_1,\ldots,E_{j-1}).
\end{equation}
无须独立即可成立，但把各环节边际成功率直接相乘需要独立假设。模型误解目标往往同时导致检索与工具错误，此时独立乘积会高估应用可靠性。

观测“执行成功”应对应业务回执；HTTP 成功只说明传输层返回了响应。用于解释动作选择的公开决策摘要可以记录结构化依据与证据标识，而无需暴露模型内部思维过程。

发布时固定整个应用制品组合，在相同任务集上比较业务结果、权限违规、人工干预、成本、恢复行为与时延。模型升级应同时检查语言任务表现、工具参数与控制器兼容性。旧长任务应继续使用兼容版本，或经过显式状态迁移；简单切换流量只影响新任务，已经提交的业务效果需要单独处理。



% BEGIN GENERATED INTERVIEWS

% END GENERATED INTERVIEWS
\chapterreferences
\myendchapterdocument
